%% kapa: exam=1 round=19 year=2008 period=2 form=A subject=economics subject_ko=경제원론 seq=38 number=38
%% source: Q-Net 감정평가사 기출문제 (cst00309) artlSeq=1001962 file=extracted/1008/1001962/2008년 제19회 감정평가사 1차시험 기출문제/2교시.hwp
%% conversion: kordoc 4.24.51 -> tools/gen1.py (auto); review=proofread-sonnet; flags=none
\begin{question}{38}
솔로우(Solow) 경제성장모형에서 생산함수가 다음과 같이 바뀌었을때 경제의 성장에 대한 다음 설명 중 적절한 것을 모두 고르면? (단, $Y$ 는 생산량, $A$ 는 생산성 수준, $K$ 는 자본을 나타낸다.)
\begin{center}$Y = AK$\end{center}
\begin{framedtext}
ㄱ. 자본에 대한 수확체감이 성립하지 않는다. \par
ㄴ. 총자본스톡 증가율은 $sA- \delta$이다. ($s$는 저축률, $\delta$는 감가상각률) \par
ㄷ. 총자본스톡 증가율이 인구증가율보다 크면 1인당 소득은 영구히 증가한다. \par
\end{framedtext}
\begin{choices}
  \item ㄱ
  \item ㄴ
  \item ㄱ, ㄷ
  \item ㄴ, ㄷ
  \item ㄱ, ㄴ, ㄷ
\end{choices}
\end{question}
\answer{5}
